\chapter{Planificación, presupuesto e impacto}
\label{chap:gestion}

\section{Planificación}

La planificación real se ha reconstruido a partir del historial del repositorio git: 35
\emph{commits} entre el 4 de marzo y el 14 de septiembre de 2026, concentrados en ocho semanas
(figura~\ref{fig:actividad-git}). El historial sólo refleja el trabajo sobre código y diseño
electrónico. Los periodos sin \emph{commits} corresponden a otras actividades (montaje mecánico,
pruebas sin cambios de código, periodo lectivo)
\completar{describir las actividades de abril a junio y de julio a agosto}.

\begin{figure}[H]
\centering
\includegraphics[width=\textwidth]{actividad_git.pdf}
\caption{\emph{Commits} por semana en el repositorio del proyecto.}
\label{fig:actividad-git}
\end{figure}

\subsection{Fases}

\begin{description}[style=nextline,leftmargin=1.5em]
  \item[F1. Exploración del accionamiento (2--10 de marzo)] Primeras pruebas de comunicación Modbus
        con una placa XIAO ESP32-S3: escaneo de direcciones, modos par y velocidad, monitor de
        terminal en Python, parámetros del accionamiento y resistencia de frenado. Culmina en un
        firmware v1.0 con un regulador de fuerza en modo par.
  \item[F2. Diseño de la placa adaptadora (12--17 de marzo)] Esquema y colocación de la
        revisión~A, con 18 copias de seguridad automáticas en seis días.
  \item[F3. Migración a Raspberry Pi Pico 2 y validación de la placa (20--30 de junio)] Programas de
        prueba de entradas y salidas, célula de carga simple y doble, RS-485 (bucle local,
        transmisión y recepción) y escaneo Modbus. Versión de doble núcleo.
  \item[F4. Firmware de posición en modo velocidad (4 de julio)] Movimiento A$\rightarrow$B con
        integración de la posición y monitor asociado.
  \item[F5. Firmware v2 y supervisión web (2--6 de septiembre)] \emph{Homing} con dos finales de
        carrera, servidor FastAPI, página de supervisión, segunda célula con ganancia, ciclos
        automáticos y grabación de ensayos.
  \item[F6. Campaña de ensayos, análisis y memoria (14--21 de septiembre)] Quince ensayos, visor de
        ensayos con modelo teórico, análisis de datos y redacción.
\end{description}

\begin{figure}[H]
\centering
\resizebox{\textwidth}{!}{%
\begin{ganttchart}[
    x unit=0.37cm, y unit chart=0.52cm, y unit title=0.55cm,
    vgrid={*1{black!8, dotted}}, hgrid,
    title label font=\scriptsize, bar label font=\scriptsize,
    milestone label font=\scriptsize\itshape,
    bar/.append style={fill=blue!45}, milestone/.append style={fill=red!70},
    bar height=0.55, group label font=\scriptsize\bfseries,
    inline=false
  ]{1}{30}
  \gantttitle{Mar}{5} \gantttitle{Abr}{4} \gantttitle{May}{4} \gantttitle{Jun}{5}
  \gantttitle{Jul}{4} \gantttitle{Ago}{5} \gantttitle{Sep}{3} \\
  \ganttbar{F1 Exploración accionamiento}{1}{2} \\
  \ganttmilestone{Firmware v1.0}{2} \\
  \ganttbar{F2 Diseño PCB rev. A}{2}{3} \\
  \ganttbar[bar/.append style={fill=black!15}]{Actividad fuera de git}{4}{15} \\
  \ganttbar{F3 Migración a Pico 2 y pruebas}{16}{18} \\
  \ganttbar{F4 Firmware de posición}{18}{18} \\
  \ganttbar[bar/.append style={fill=black!15}]{Actividad fuera de git}{19}{26} \\
  \ganttbar{F5 Firmware v2 y web}{27}{28} \\
  \ganttmilestone{Primer ciclo grabado}{28} \\
  \ganttbar{F6 Ensayos y análisis}{29}{29} \\
  \ganttbar{Redacción de la memoria}{29}{30}
\end{ganttchart}%
}
\caption{Diagrama de Gantt reconstruido (semanas desde el 2 de marzo de 2026).}
\label{fig:gantt}
\end{figure}

\section{Presupuesto}

\subsection{Recursos humanos}

La dedicación de un TFG se fija en créditos ECTS, cada uno de 25 a 30 horas de trabajo
\cite{rd1125}. La tabla~\ref{tab:horas} reparte una dedicación estimada por fases
\verificar{ajustar las horas y los créditos del TFG en la titulación}. Para valorar el trabajo se
toma como supuesto un coste de 30~€/h para un ingeniero o ingeniera junior y de
45~€/h para la tutoría.

\begin{table}[H]
\centering
\caption{Dedicación estimada y coste de personal.}
\label{tab:horas}
\begin{tabular}{@{}lrrr@{}}
\toprule
Actividad & Horas & €/h & Coste (€) \\
\midrule
Estudio y documentación & 30 & 30 & 900 \\
Exploración del accionamiento (F1) & 35 & 30 & 1.050 \\
Diseño de la placa (F2) & 45 & 30 & 1.350 \\
Pruebas de placa y comunicaciones (F3) & 45 & 30 & 1.350 \\
Firmware (F4--F5) & 70 & 30 & 2.100 \\
Software de supervisión (F5) & 50 & 30 & 1.500 \\
Ensayos y análisis (F6) & 25 & 30 & 750 \\
Redacción de la memoria & 50 & 30 & 1.500 \\
\midrule
Subtotal trabajo técnico & 350 & & 10.500 \\
Tutoría & 25 & 45 & 1.125 \\
\midrule
\textbf{Total recursos humanos} & & & \textbf{11.625} \\
\bottomrule
\end{tabular}
\end{table}

\subsection{Material}

\begin{longtable}{@{}L{0.42\textwidth}rrr@{}}
\caption{Material del sistema.}\label{tab:material}\\
\toprule
Elemento & Cant. & Precio unit. (€) & Total (€) \\
\midrule
\endfirsthead
\toprule
Elemento & Cant. & Precio unit. (€) & Total (€) \\
\midrule
\endhead
Raspberry Pi Pico 2 & 1 & \completar{} & \\
Transceptor RS-485 THVD1400 & 1 & \completar{} & \\
Interruptor de lado alto TPS4H000-Q1 & 1 & \completar{} & \\
Amplificador operacional MCP6294 & 1 & \completar{} & \\
Diodos zéner BZX384 & 4 & \completar{} & \\
Resistencias y condensadores SMD 0805 & 25 & \completar{} & \\
Bornas de tornillo y tiras de pines & -- & \completar{} & \\
Fabricación de la PCB & 1 & \completar{} & \\
Load Cell 2 Click (ZSC31014) & 1 & \completar{} & \\
Célula de carga & 1 & \completar{} & \\
Servomotor paso a paso de lazo cerrado con accionamiento & 1 & \completar{} & \\
Fuente de alimentación del accionamiento & 1 & \completar{} & \\
Finales de carrera & 2 & \completar{} & \\
Caracola, barra, cable, estructura y masa de \SI{2}{\kilo\gram} & -- & \completar{} & \\
\midrule
\textbf{Total material} & & & \completar{} \\
\bottomrule
\end{longtable}

\subsection{Equipos y software}

Todo el software empleado es libre o gratuito (KiCad, Arduino IDE con el núcleo Arduino-Pico,
Python, FastAPI, Chart.js, git, \LaTeX), por lo que su coste es nulo. Los equipos se imputan por
amortización lineal,
\begin{equation}
  A = C\,\frac{t_{\text{uso}}}{t_{\text{vida}}},
\end{equation}
con $C$ el coste de adquisición, $t_{\text{uso}}$ el tiempo de uso en el proyecto y $t_{\text{vida}}$ la
vida útil (tabla~\ref{tab:amortizacion}).

\begin{table}[H]
\centering
\caption{Amortización de equipos.}
\label{tab:amortizacion}
\begin{tabular}{@{}lrrrr@{}}
\toprule
Equipo & Coste (€) & Vida útil (meses) & Uso (meses) & Amortización (€) \\
\midrule
Ordenador portátil & \completar{} & 48 & 7 & \\
Polímetro & \completar{} & 60 & 7 & \\
Fuente de laboratorio & \completar{} & 60 & 7 & \\
\midrule
\textbf{Total} & & & & \completar{} \\
\bottomrule
\end{tabular}
\end{table}

\subsection{Resumen}

\begin{table}[H]
\centering
\caption{Presupuesto total.}
\label{tab:presupuesto}
\begin{tabular}{@{}lr@{}}
\toprule
Concepto & Importe (€) \\
\midrule
Recursos humanos & 11.625 \\
Material & \completar{} \\
Amortización de equipos & \completar{} \\
\midrule
Subtotal & \completar{} \\
Costes indirectos (\SI{10}{\percent}, supuesto) & \completar{} \\
\midrule
Base imponible & \completar{} \\
IVA (\SI{21}{\percent}) & \completar{} \\
\midrule
\textbf{Total} & \completar{} \\
\bottomrule
\end{tabular}
\end{table}

\section{Impacto}

\subsection{Impacto técnico y económico}

El resultado del capítulo~\ref{chap:modelo} tiene una consecuencia directa en el dimensionado:
para un recorrido dado, un perfil de par constante reduce el par de pico que debe dar el motor. Con
la geometría nominal la reducción es del \SI{24.7}{\percent}, lo que puede permitir un motor y un
accionamiento de un tamaño inferior, con menor coste y menor consumo de materiales. Además, a
velocidad constante las pérdidas por efecto Joule en el motor son proporcionales a
$\int\tau^2\,\mathrm{d}\varphi$. Para un trabajo fijo, esa integral es mínima cuando el par es
constante (desigualdad de Cauchy-Schwarz). Con la geometría nominal, el perfil ideal reduce esas
pérdidas un \SI{3.5}{\percent} respecto al tambor cilíndrico.

La plataforma desarrollada es además reutilizable. La placa adaptadora, el firmware y la aplicación
web forman un banco genérico para cualquier eje accionado por Modbus con medida de fuerza, y su
coste de material es muy inferior al de un sistema de adquisición comercial.

\subsection{Objetivos de Desarrollo Sostenible}

\begin{table}[H]
\centering
\caption{Relación del trabajo con los Objetivos de Desarrollo Sostenible \cite{agenda2030}.}
\label{tab:ods}
\begin{tabularx}{\textwidth}{@{}lX@{}}
\toprule
ODS & Contribución \\
\midrule
4. Educación de calidad & Banco de ensayo abierto y documentado, útil como práctica de laboratorio de mecanismos, accionamientos y adquisición de datos. \\
7. Energía asequible y no contaminante & Reducción de las pérdidas en el motor mediante un perfil de par uniforme. \\
9. Industria, innovación e infraestructura & Integración de tecnologías industriales (Modbus, RS-485) con herramientas libres de bajo coste. \\
12. Producción y consumo responsables & Dimensionado de actuadores más pequeños; módulo microcontrolador en zócalo, reutilizable; datos brutos conservados para no repetir ensayos. \\
\bottomrule
\end{tabularx}
\end{table}

\subsection{Aspectos éticos, legales y de seguridad}

El prototipo es un equipo de laboratorio que no se comercializa. Aun así, su evolución hacia una
máquina de uso por terceros tendría que considerar la Directiva 2006/42/CE de máquinas, que el
Reglamento (UE) 2023/1230 sustituye a partir de enero de 2027, la Directiva 2014/30/UE de
compatibilidad electromagnética y la Directiva 2011/65/UE sobre sustancias peligrosas (RoHS). La
sección~\ref{sec:arq-seguridad} discute las limitaciones de seguridad del diseño actual frente a
ISO 13849-1 \cite{iso13849}.

En cuanto a la ética profesional, se ha priorizado la trazabilidad de los datos: se conservan las
medidas brutas, las figuras se generan con un programa reproducible y las limitaciones y los datos
provisionales se señalan de forma explícita en la memoria.
